% Section 7 — Results. DATA STATUS: the confirmatory campaign is being executed on the
% approved Linux/container host; this section fixes the reporting structure and is filled
% exclusively from frozen run records by tools/analyze.py and tools/make-figures.py.
% No value in this section is typed by hand; every float is marked PENDING until generated.
\section{Results}
\label{sec:results}

\textbf{Data status.} The confirmatory campaign is executed from the frozen matrix and analysis
plan of Section~\ref{sec:method}; at preparation time it is in progress on the approved
Linux/container host, and the Windows feasibility run stopped before any confirmatory cell. Every
table and figure here is generated by the frozen pipeline from accepted run records; none is typed
by hand, and each is marked \emph{PENDING} until the pipeline emits its content.

\subsection{Campaign integrity and harness validation}
\label{sec:res-integrity}

Before any confirmatory run is analysed, the eight harness gates of Section~\ref{sec:method}
must pass: the negative control must duplicate, the positive controls must not, the checker must
detect every synthetic mutation, every armed injector must leave a complete fired/restored
receipt pair, the ground-truth digests must agree between generator and oracle, and the shared
configuration hash must be identical across mechanisms. Table~\ref{tab:validation} records the
gate evidence and the per-cell exclusion of runs whose fault did not fire; no other exclusion is
permitted.

\IfFileExists{results/analysis/tables/t1-validation.tex}{\input{results/analysis/tables/t1-validation.tex}}{%
\begin{center}
\footnotesize
\fbox{\parbox{0.92\linewidth}{\centering\textbf{Table PENDING.} Gates G1--G8: evidence and
per-cell exclusion rates, generated from frozen validation records.}}
\captionof{table}{Harness-validation evidence and exclusions (generated from frozen run records;
PENDING until the campaign completes).}
\label{tab:validation}
\end{center}}

\subsection{RQ1: where exactly-once ends}
\label{sec:res-rq1}

Table~\ref{tab:correctness} will report, per mechanism and fault condition, violating runs and
the one-sided Clopper--Pearson 95\,\% upper bound (rule of three for zero-event cells) under the
frozen decision rules. The pre-registered prediction is that M0 violates the one-entry-per-event
invariant (I3) whenever the duplicate window is reachable, M1-kafka violates no ledger invariant,
and M1x violates in the classes the crash-point model maps to the external commit
(C3/C4/C9/C10); the early-ordering probe is predicted to show the loss violation (I4) for every
two-commit mechanism without database-resident progress. Table~\ref{tab:counts} decomposes
violations into duplicate, lost, imbalanced and mutated postings, and the report will state which
invariants were never observed to fail --- with exact upper bounds, never treating zero as a
proof. Figure~\ref{fig:correctness} places the outcomes against their measured cost.

\IfFileExists{results/analysis/tables/t2-correctness.tex}{\input{results/analysis/tables/t2-correctness.tex}}{%
\begin{center}
\footnotesize
\fbox{\parbox{0.92\linewidth}{\centering\textbf{Table PENDING.} Correctness by mechanism
$\times$ fault class: violating runs, $N$, one-sided 95\,\% upper bound.}}
\captionof{table}{Correctness by mechanism and fault class (generated from frozen run records;
PENDING until the campaign completes).}
\label{tab:correctness}
\end{center}}

\IfFileExists{results/analysis/tables/t3-counts.tex}{\input{results/analysis/tables/t3-counts.tex}}{%
\begin{center}
\footnotesize
\fbox{\parbox{0.92\linewidth}{\centering\textbf{Table PENDING.} Duplicate, lost, imbalanced and
mutated posting counts by mechanism.}}
\captionof{table}{Violation counts by invariant and mechanism (generated from frozen run records;
PENDING until the campaign completes).}
\label{tab:counts}
\end{center}}

\begin{center}
\IfFileExists{figures/fig-correctness.pdf}{\includegraphics[width=\linewidth]{figures/fig-correctness}}{%
\fbox{\parbox{0.92\linewidth}{\centering\textbf{Figure PENDING.} Panel~A: per-run violation
proportion by mechanism and fault class with 95\,\% intervals. Panel~B: failure classes closed
versus p99 latency overhead. Generated from frozen run records.}}}
\captionof{figure}{Correctness evidence. Panel~A: per-run ledger-invariant violation proportion by
mechanism and fault class (Clopper--Pearson 95\,\% upper bounds; structurally unreachable cells
are n/a and never counted as passes). Panel~B: failure classes whose invariant is preserved
versus p99 latency overhead relative to the naive baseline; higher and further left is better.
PENDING until the frozen analysis completes.}
\label{fig:correctness}
\end{center}

\subsection{RQ2: the cost of closing the boundary}
\label{sec:res-rq2}

Table~\ref{tab:performance} will report throughput, end-to-end latency (p50/p99/p99.9) and
time-to-consistent-state by mechanism, load and partition count, with median/IQR and bootstrap
95\,\% intervals; the pre-registered contrasts (Friedman/Nemenyi, Cliff's delta) are stated
without calling a mechanism ``fast'' or ``slow'' outside the frozen cells.
Figure~\ref{fig:performance} shows latency ECDFs (panel~A) and throughput by load and partitions
(panel~B). Closing the duplicate window has a cost: M3 maintains an index, M4 moves progress into
the posting transaction, and M2 adds an inbox claim, outbox row and relay thread.

\IfFileExists{results/analysis/tables/t4-performance.tex}{\input{results/analysis/tables/t4-performance.tex}}{%
\begin{center}
\footnotesize
\fbox{\parbox{0.92\linewidth}{\centering\textbf{Table PENDING.} Performance by mechanism
$\times$ load $\times$ partitions: throughput, p50/p99/p99.9, time-to-consistency.}}
\captionof{table}{Performance by mechanism, load and partition count (generated from frozen run
records; PENDING until the campaign completes).}
\label{tab:performance}
\end{center}}

\begin{center}
\IfFileExists{figures/fig-latency-ecdf.pdf}{\includegraphics[width=\linewidth]{figures/fig-latency-ecdf}}{%
\fbox{\parbox{0.92\linewidth}{\centering\textbf{Figure PENDING.} Panel~A: end-to-end latency
ECDFs per mechanism (fault-free runs). Panel~B: committed-event throughput by load and
partition count. Generated from frozen run records.}}}
\captionof{figure}{Performance evidence. Panel~A: end-to-end latency ECDFs per mechanism for fault-free
runs (open-loop load; warm-up excluded). Panel~B: committed-event throughput by load and
source-partition count (median across repetitions, IQR shown). PENDING until the frozen analysis
completes.}
\label{fig:performance}
\end{center}

\subsection{RQ3: minimal sufficient mechanism sets}
\label{sec:res-rq3}

Section~\ref{sec:decision} combines the frozen correctness outcomes with the measured cost
contrasts under the pre-registered derivation rule: the duplicate window is closed by M3's index
or M4's atomic progress; the loss window only by a mechanism making progress atomic with the
effect (M4, M2, M4+M3, M2+M3), never by committing progress ahead of a durable effect (which the
`early' probe violates); M1x closes neither; M1x+M3 closes duplicates while retaining
Kafka-internal exactly-once for publication. Measured cost decides among sufficient sets.

\subsection{Hypothesis verdicts}
\label{sec:res-hypotheses}

The pre-registered decision for each hypothesis is computed by the frozen decision rules of
Section~\ref{sec:method} from the run records, and a verdict is recorded exactly as the rule
computes it (confirmed, refuted, or inconclusive). Until the campaign completes every verdict
reads \emph{pending}; the values are populated by the frozen analysis pipeline, not by hand.

The five pre-registered hypotheses are: \textbf{H1} M1x violates and M1-kafka does not;
\textbf{H2} M2 is sufficient but costs throughput/latency vs.\ M3/M4; \textbf{H3} M3 alone yields
zero duplicates but losses under early ordering; \textbf{H4} M4 or M4+M3 is sufficient with lower
latency than M2; \textbf{H5} no mechanism dominates. Their verdicts are computed by the frozen
decision rules:
H1 \ResHypothesisOneVerdict; H2 \ResHypothesisTwoVerdict; H3 \ResHypothesisThreeVerdict;
H4 \ResHypothesisFourVerdict; H5 \ResHypothesisFiveVerdict. The cells read \emph{pending} until
the analysis pipeline runs, and a verdict is recorded exactly as the rule computes it (confirmed,
refuted or inconclusive).
